\subsection{\texorpdfstring{第~\ref{sec:manifolds-of-constant-curvature}~节}{常曲率流形}}\label{subsec:solutions-manifolds-of-constant-curvature}

\begin{description}[font=\normalfont,leftmargin=2em,labelwidth=1.5em]
\RGsolution{ex:discrete-isometry-groups-of-plane}{discrete-isometry-groups-of-plane}
\begin{enumerate}[label=(\alph*),ref=\alph*,leftmargin=*]
\item $\mathbb{R}^2$ 的等距映射群是
\[
  \operatorname{Iso}(\mathbb{R}^2)=\{Ax+b\mid A\in O(2),\ b\in\mathbb{R}^2\}.
\]
将 $\mathbb{R}^2$ 与复平面 $\mathbb{C}$ 等同, 可以将任意保持定向的元素 $f\in\operatorname{Iso}(\mathbb{R}^2)$ 写为 $f(z)=e^{i\theta}z+b$, 其中 $\theta\in\mathbb{R}$ 且 $b\in\mathbb{C}$. $f$ 的不动点由
\[
  f(z)=z\iff(1-e^{i\theta})z=b
\]
给出. 因而 $f$ 没有不动点当且仅当 $e^{i\theta}=1$, 即 $f$ 是平移. 另一方面, 任意反转定向的元素 $g\in\operatorname{Iso}(\mathbb{R}^2)$ 均可写为 $g(z)=e^{i\theta}\bar z+b$, 其中 $\theta\in\mathbb{R}$ 且 $b\in\mathbb{C}$. 例如, 沿实轴的滑移反射由 $h(z)=\bar z+\xi$ 给出 (其中 $\xi\in\mathbb{R}$). 设 $r(z)=e^{i\frac{\theta}{2}}z$ 为转角 $\frac{\theta}{2}$ 的旋转, $t(z)=z+i\eta$ (其中 $\eta\in\mathbb{R}$) 为沿 $i\eta\in i\mathbb{R}$ 的平移. 则
\[
  (r\circ h\circ t\circ r^{-1})(z)
  =e^{i\frac{\theta}{2}}\left(\overline{e^{-i\frac{\theta}{2}}z+i\eta}+\xi\right)
  =e^{i\theta}\bar z+e^{i\frac{\theta}{2}}(\xi-i\eta)
\]
是一个滑移反射 (轴的斜率为 $\tan\frac{\theta}{2}$, 且它到原点的距离为 $|\eta|$). 当 $\xi,\eta$ 满足
\[
  \xi+i\eta=e^{i\frac{\theta}{2}}\bar b
\]
时, 这个映射等于 $g$, 所以 $\mathbb{R}^2$ 上任意反转定向的等距映射都是滑移反射. 只要沿反射轴的平移不为零, 即 $e^{-i\frac{\theta}{2}}b\notin i\mathbb{R}$, 这些映射显然没有不动点.
\item 设 $\Gamma$ 是固有且自由地作用于 $\mathbb{R}^2$ 的 $\operatorname{Iso}(\mathbb{R}^2)$ 的离散子群. 那么 $\Gamma$ 只能包含平移和滑移反射.

先设 $\Gamma$ 只包含平移, 设 $t_1\in\Gamma$ 是使 $t_1(0)$ 的长度最小的 $\Gamma$ 中平移 (因 $\Gamma$ 的作用为固有而存在). 那么由 $t_1$ 生成的群 $\Gamma_1:=\langle t_1\rangle\subset\Gamma$ 包含 $\Gamma$ 中与 $t_1$ 同方向的所有平移. 实际上, 若 $t$ 是 $\Gamma\setminus\Gamma_1$ 中与 $t_1$ 同方向的平移, 则取 $m\in\mathbb{Z}$ 使 $t_1^m$ 是 $\Gamma_1$ 中使 $t_1^m(0)$ 最接近 $t(0)$ 的元素, 这时 $t^{-1}t_1^m(0)$ 的长度小于 $t_1(0)$ 的长度, 与初始假设矛盾.

若 $\Gamma_1\ne\Gamma$, 设 $t_2\in\Gamma\setminus\Gamma_1$ 使 $t_2(0)$ 的长度最小. 则 $t_1$ 与 $t_2$ 在 $\mathbb{R}^2$ 中生成一个格, 且 $\Gamma=\langle t_1,t_2\rangle$. 实际上, 若在 $\Gamma\setminus\langle t_1,t_2\rangle$ 中存在元素 $t$, 则取 $m,n\in\mathbb{Z}$ 使 $t_1^m t_2^n(0)$ 最接近 $t(0)$, 这时 $t^{-1}t_1^m t_2^n(0)$ 的长度或者小于 $t_1(0)$ 的长度, 或者小于 $t_2(0)$ 的长度.

若 $\Gamma$ 含有滑移反射, 设 $g\in\Gamma$ 为使 $g^2(0)$ 的长度最小的滑移反射 (因 $\Gamma$ 的作用为固有而存在). 若 $\Gamma_1:=\langle g\rangle\ne\Gamma$, 则 $\Gamma\setminus\Gamma_1$ 含有平移 (若 $g_1\in\Gamma\setminus\Gamma_1$ 为滑移反射, 则 $g^{-1}g_1$ 是 $\Gamma\setminus\Gamma_1$ 中的平移). 设 $t\in\Gamma\setminus\Gamma_1$ 是使 $t(0)$ 的长度最小的平移 (因 $\Gamma$ 的作用为固有而存在). 适当选取坐标后可设 $g(z)=\bar z+\xi$, 其中 $\xi\in\mathbb{R}$. 设 $t(z)=z+b$, 其中 $b=\alpha+i\beta\in\mathbb{C}$ ($\alpha,\beta\in\mathbb{R}$). 则
\[
  (g^{-1}\circ t\circ g)(z)=g^{-1}(\bar z+\xi+b)=z+\bar b,
\]
从而 $\Gamma$ 含有沿 $b$ 与 $\bar b$ 的平移, 故也含有沿 $2\alpha$ 和 $2i\beta$ 的平移. 因为 $g^2(0)=2\xi$ 的长度最小, 容易检验 $2\alpha$ 必为 $\xi$ 的整数倍, 否则便可构造满足 $|\widetilde g^2(0)|<|g^2(0)|$ 的滑移反射 $\widetilde g$. 由于 $t(0)$ 的长度最小, 可知 $\alpha$ 只能为 $0$, $\pm\frac{\xi}{2}$ 或 $\pm\xi$; 但是 $\alpha$ 不能是 $\pm\frac{\xi}{2}$, 因为在这种情形下 $\Gamma$ 将含有反射 $z\mapsto\bar z$, 也不能是 $\pm\xi$, 因为在这种情形下 $\Gamma$ 将含有反射 $z\mapsto\bar z+i\beta$. 因此, $t$ 是沿 $i\beta$ 的平移, 从而 $0\in\mathbb{R}^2$ 在 $\langle g,t\rangle$ 下的轨道与它在由沿 $\xi$ 和 $i\beta$ 的平移所生成的群下的轨道相同. 与上面类似的论证表明 $\Gamma=\langle g,t\rangle$.
\end{enumerate}

\RGsolution{ex:mobius-band-and-klein-bottle-quotients}{mobius-band-and-klein-bottle-quotients}
为确定 $\mathbb{R}^2/\langle f\rangle$, 首先注意开半平面 $\{(x,y)\in\mathbb{R}^2\mid x<0\}$ 中每一点都等价于开半平面 $\{(x,y)\in\mathbb{R}^2\mid x>0\}$ 中的一点. 另一方面,
\[
  f^2(x,y)=f(f(x,y))=f(-x,y+1)=(x,y+2).
\]
所以 $\mathbb{R}^2$ 中任一点都等价于条带
\[
  S=\{(x,y)\in\mathbb{R}^2\mid x\ge0\text{ 且 }0\le y\le2\}
\]
中的一点. 因而 $\mathbb{R}^2/\langle f\rangle$ 与 $S/{\sim}$ 同胚, 其中等价关系 $\sim$ 在 $S$ 上由
\[
  (x,0)\sim(x,2)\quad\text{和}\quad(0,y)\sim(0,y+1)
\]
定义. 现在, $S/{\sim}$ 显然与半无限圆柱面
\[
  C=\{(x,y,z)\in\mathbb{R}^3\mid x^2+y^2=1\text{ 且 }z\ge0\}
\]
在边界上按 $(x,y,0)\sim(-x,-y,0)$ 等同所得的商空间同胚 (参见 Fig.~\ref{fig:mobius-band-as-cylinder-quotient}). 这个商空间又显然与去掉一个闭圆盘的射影平面同胚, 后者与 Möbius 带 (去掉边界) 同胚.
\RGfigure[.59\textwidth]{mobius-band-as-cylinder-quotient}{$\mathbb{R}^2/\langle f\rangle$}

为确定 $\mathbb{R}^2/\langle f,g\rangle$, 注意 $\mathbb{R}^2$ 中任一点都等价于矩形 $R=[0,1]\times[0,2]$ 中的一点. 此外, 因为
\[
  g(f(x,y))=g(-x,y+1)=(1-x,y+1),
\]
可见 $\mathbb{R}^2$ 中任一点实际上都等价于正方形 $Q=[0,1]\times[0,1]$ 中的一点. 因此 $\mathbb{R}^2/\langle f,g\rangle$ 与 $Q/{\sim}$ 同胚, 其中等价关系 $\sim$ 在 $Q$ 上由
\[
  (0,y)\sim(1,y)\quad\text{和}\quad(x,0)\sim(1-x,1)
\]
定义. 这恰好是一个 Klein 瓶.

\RGsolution{ex:orientation-preserving-isometries-of-hyperbolic-plane}{orientation-preserving-isometries-of-hyperbolic-plane}
\begin{enumerate}[label=(\alph*),ref=\alph*,leftmargin=*]
\item 首先注意映射
\[
  f(z):=\frac{az+b}{cz+d}
\]
满足 $f'(z)=\frac1{(cz+d)^2}$ 和 $\operatorname{Im}(f(z))=\frac{y}{|cz+d|^2}$, 其中 $z=x+iy$. 容易验证 $g$ 将 $H^2$ 映到自身上, 且若 $f(z)=v(x,y)+iv(x,y)$, 则
\begin{align*}
  (f^*g)_z\left(\frac{\partial}{\partial x},\frac{\partial}{\partial x}\right)
  &=g_{f(z)}\left(f_*\frac{\partial}{\partial x},f_*\frac{\partial}{\partial x}\right)\\
  &=g_{f(z)}\left(\frac{\partial v}{\partial x}\frac{\partial}{\partial x}
    +\frac{\partial v}{\partial x}\frac{\partial}{\partial y},
    \frac{\partial v}{\partial x}\frac{\partial}{\partial x}
    +\frac{\partial v}{\partial x}\frac{\partial}{\partial y}\right)\\
  &=g_{f(z)}\left(\frac{\partial v}{\partial x}\frac{\partial}{\partial x}
    -\frac{\partial v}{\partial y}\frac{\partial}{\partial y},
    \frac{\partial v}{\partial x}\frac{\partial}{\partial x}
    -\frac{\partial v}{\partial y}\frac{\partial}{\partial y}\right)\\
  &=\frac1{v^2}\left[\left(\frac{\partial v}{\partial x}\right)^2+
    \left(\frac{\partial v}{\partial y}\right)^2\right]=\frac1{v^2}|f'(z)|^2\\
  &=\frac1{y^2}=g_p\left(\frac{\partial}{\partial x},\frac{\partial}{\partial x}\right),
\end{align*}
这里使用了 $f$ 的 Cauchy–Riemann 方程 (例如见 \RGcite{Ahl79}). 类似地, 可知
\[
  (f^*g)_p\left(\frac{\partial}{\partial x},\frac{\partial}{\partial y}\right)=0
  \quad\text{且}\quad
  (f^*g)_p\left(\frac{\partial}{\partial y},\frac{\partial}{\partial y}\right)=\frac1{y^2},
\]
所以 $f$ 是 $H^2$ 的等距映射. 此外, $\det(\mathrm{d}f)=|f'(z)|^2>0$, 所以 $f$ 保持定向.
\item\label{sol:psl-transitive-on-unit-tangents} 回顾, 等距映射将测地线送到测地线. 所以只需证明: 给定 $H^2$ 中两点 $p,q$ 和分别在 $p,q$ 处的单位向量 $v,w$, 存在 $g\in PSL(2,\mathbb{R})$ 使 $g(p)=q$ 和 $(\mathrm{d}g)_pv=w$ (这里也用 $g$ 表示由 $g(z)=g\cdot z$ 给出的映射). 为此, 首先看到 $i$ 的轨道是整个 $H^2$. 实际上, 对任意 $z_0=x_0+iy_0$, 其中 $y_0>0$, 映射 $f_{z_0}(z):=y_0z+x_0$ 对应于矩阵
\[
  \begin{pmatrix}\sqrt{y_0}&\dfrac{x_0}{\sqrt{y_0}}\\[2pt]
                  0&\dfrac1{\sqrt{y_0}}\end{pmatrix}\in PSL(2,\mathbb{R}),
\]
并把 $i$ 送到 $z_0$. 此外, 与矩阵
\[
  \begin{pmatrix}\cos\theta&\sin\theta\\-\sin\theta&\cos\theta\end{pmatrix}\in PSL(2,\mathbb{R})
\]
对应的映射 $r_\theta:H^2\to H^2$ 固定 $i$ (即对每个 $\theta$ 有 $r_\theta(i)=i$), 且在 $i$ 处的向量上传递地作用, 因为
\[
  (\mathrm{d}r_\theta)_i=
  \begin{pmatrix}\cos(2\theta)&-\sin(2\theta)\\\sin(2\theta)&\cos(2\theta)\end{pmatrix}
\]
(它们在 $T_iH^2$ 上如 $SO(2)$ 一样作用). 因此映射 $g=f_q\circ r_\theta\circ f_p^{-1}$, 其中 $\theta$ 满足 $(\mathrm{d}r_\theta)_i(\mathrm{d}f_p^{-1})_pv=(\mathrm{d}f_q^{-1})_qw$, 把 $p$ 送到 $q$, 且使得 $(\mathrm{d}g)_pv=w$.
\item\label{sol:psl-preserves-geodesic-pairs} 这直接由 (\ref{sol:psl-transitive-on-unit-tangents}) 得出.
\item 设 $f:H^2\to H^2$ 为保持定向的等距映射. 若 $f$ 有两个不动点, 则它必固定经过这两点的测地线. 设 $z$ 是这条测地线上的一点, 选取正定向标准正交标架 $\{E_1,E_2\}\in T_zH^2$, 其中 $E_1$ 与此测地线相切. 则 $(\mathrm{d}f)_z E_1=E_1$, 并且因为 $f$ 是保持定向的等距映射, 有 $(\mathrm{d}f)_zE_2=E_2$. 所以 $(\mathrm{d}f)_z$ 必为恒等映射, 从而 $f$ 固定所有经过 $z$ 的测地线, 即 $f$ 为恒等映射.

现设 $f:H^2\to H^2$ 为任意保持定向的等距映射, 取两点 $z_1,z_2\in H^2$, 并设 $g$ 为 (\ref{sol:psl-preserves-geodesic-pairs}) 所给的使 $g(z_1)=f(z_1)$ 和 $g(z_2)=f(z_2)$ 的 $PSL(2,\mathbb{R})$ 元素. 那么 $f\circ g^{-1}$ 是有两个不动点的保持定向的等距映射, 因此它必为恒等映射, 即 $f=g$.
\end{enumerate}
\end{description}
